\section{From endpoint costs to metric medians}\label{sec:median}

We now convert the optimal representations into a geometric property
of the polar body. A \emph{metric median} of three points in a normed
space is a point that lies between every pair: if the points are
$x_0,x_1,x_2$, a median $m$ satisfies
\[
 \|x_i-x_j\|=\|x_i-m\|+\|m-x_j\|
 \qquad(0\le i<j\le2).
\]
The median need not be unique. Its existence will imply the ball
intersection property that characterizes linear Hanner bodies.

\begin{proposition}[Medians of an equality body]\label{prop:metric-median}
Let $K\subset\R^n$, $n\ge1$, be an origin-symmetric convex body with
$P(K)=4^n/n!$, and put $Q=K^\circ$. Every triple of points in
$(\R^n,p_Q)$ has a metric median.
\end{proposition}

\begin{proof}
Write $\|\cdot\|=p_Q$. Translation reduces the triple to $0,x,y$.
If two points coincide, that repeated point is a median, so assume
that $0,x,y$ are distinct. Set
\[
 p_1=\|x\|,\qquad p_2=\|y\|,\qquad p=\|x-y\|,
 \qquad s=p_1+p_2,\qquad a=p_1-p_2.
\]
All three distances are positive. The triangle and reverse triangle
inequalities give $s\ge p\ge|a|$. Define
\begin{equation}\label{eq:median-radii}
 A_+=\frac{p+a}{2},\qquad A_-=\frac{p-a}{2},
 \qquad c=\frac{s-p}{2}.
\end{equation}
These numbers are nonnegative, with
$A_++A_-=p$, $c+A_+=p_1$, and $c+A_-=p_2$.

\paragraph{The compact set of aggregate directions.}
Let
\begin{equation}\label{eq:aggregate-set}
 D_0=\{r_+-r_-:\ r_++r_-=x-y,\quad
                  \|r_+\|\le A_+,\quad\|r_-\|\le A_-\}.
\end{equation}
Our goal is to prove $x+y\in D_0+(s-p)Q$. Such membership supplies
$r_+,r_-$ for which $m=x-r_+=y+r_-$ has distances to $0,x,y$
bounded by $c,A_+,A_-$, respectively; the triangle inequalities then
force the three median equalities. We shall prove this membership by
estimating support functions in each direction.

The set $D_0$ is compact and convex: it is the linear image of the intersection
of $A_+Q\times A_-Q$ with the affine subspace $r_++r_-=x-y$. It is nonempty, since
$r_\pm=(A_\pm/p)(x-y)$ satisfy the constraints. This also covers
$A_+=0$ or $A_-=0$.

\paragraph{Endpoint comparison.}
Fix $u\in\operatorname{int}K$ and let
$X_\delta=(\delta u,1-\delta)$, $0<\delta<1$. The lifted body
$B=K\oplus_1[-1,1]$ contains $X_\delta$ in its interior because
\[
 p_B(X_\delta)=\delta p_K(u)+1-\delta<1.
\]
Apply Proposition~\ref{prop:optimal-representation} with this $X_\delta$
and $(q,a)=(x-y,p_1-p_2)$. Its total weight is
$T=\max\{\|q\|,|a|\}=p$. We obtain rows
$(b_i,\sigma_i)\in Q\times\{-1,1\}$ and weights $\alpha_i\ge0$
satisfying \eqref{eq:optimal-representation} and
\eqref{eq:optimal-cost}. Their aggregates
\[
 r_\pm=\sum_{\sigma_i=\pm1}\alpha_i b_i
\]
have $r_++r_-=x-y$. Since
$\sum_{\sigma_i=\pm1}\alpha_i=(p\pm a)/2=A_\pm$, they satisfy
$\|r_\pm\|\le A_\pm$ and hence $r_+-r_-\in D_0$.
No choice of these representations across different $u$ or $\delta$
is required.

For $b\in Q$ and $\sigma=\pm1$, evenness of $\lambda$ gives
\[
 \lambda\bigl(\langle(b,\sigma),X_\delta\rangle\bigr)
 =\lambda\bigl(1-\delta(1-\sigma\langle b,u\rangle)\bigr).
\]
The factors $1-\sigma\langle b,u\rangle$ belong to
$[1-p_K(u),1+p_K(u)]$, a compact subinterval of $(0,\infty)$.
Lemma~\ref{lem:lens-endpoint} therefore gives an error
$\varepsilon_u(\delta)\to0$ such that
\begin{equation}\label{eq:median-endpoint-error}
 \left|
 \frac{\lambda(\langle(b,\sigma),X_\delta\rangle)}
      {\lambda(1-\delta)}
       -1+\sigma\langle b,u\rangle
 \right|\le\varepsilon_u(\delta)
 \quad(b\in Q,\ \sigma=\pm1).
\end{equation}
In particular, the normalized cost of the optimal representation is
at least
\[
 p-\langle r_+-r_-,u\rangle-p\varepsilon_u(\delta)
 \ge p-h_{D_0}(u)-p\varepsilon_u(\delta).
\]
Compare this cost, using \eqref{eq:optimal-cost}, with the two-term
representation
\[
 (x-y,a)=p_1(x/p_1,1)+p_2(-y/p_2,-1).
\]
Both rows belong to $Q\times\{-1,1\}$. By
\eqref{eq:median-endpoint-error}, their normalized cost tends to
$s-\langle x+y,u\rangle$. Letting $\delta\downarrow0$ yields
\begin{equation}\label{eq:median-support-interior}
 \langle x+y,u\rangle\le h_{D_0}(u)+s-p
 \qquad(u\in\operatorname{int}K).
\end{equation}
The error was uniform only over $b,\sigma$ for the fixed $u$; that
is all this limit requires.

\paragraph{Separation and the median.}
Continuity extends \eqref{eq:median-support-interior} to $u\in K$.
For any nonzero $v\in\R^n$, the point
$u=v/h_Q(v)$ lies on $\partial K$, by \eqref{eq:gauge-support}.
Homogeneity of support functions gives
\[
 \langle x+y,v\rangle\le h_{D_0}(v)+(s-p)h_Q(v).
\]
The same inequality holds at $v=0$. The right side is the support
function of the compact convex set $D_0+(s-p)Q$, with the convention
$0Q=\{0\}$. Convex separation therefore implies
\begin{equation}\label{eq:median-membership}
 x+y\in D_0+(s-p)Q.
\end{equation}
Choose the corresponding pair $r_+,r_-$ in
\eqref{eq:aggregate-set}, and put
\[
 m=x-r_+=y+r_-.
\]
Then $2m=x+y-(r_+-r_-)$, so \eqref{eq:median-membership} gives
\[
 \|m\|\le c,\qquad
 \|x-m\|\le A_+,\qquad
 \|y-m\|\le A_-.
\]
The triangle inequality forces all three bounds to be equalities:
\[
 \begin{aligned}
 p_1&\le\|m\|+\|x-m\|\le c+A_+=p_1,\\
 p_2&\le\|m\|+\|y-m\|\le c+A_-=p_2.
 \end{aligned}
\]
Thus $m$ lies between $0,x$ and between $0,y$, and
\[
 \|x-m\|+\|m-y\|=A_++A_-=p=\|x-y\|.
\]
This proves the median property. The argument permits $c=0$ or either
$A_\pm=0$, so saturated triangle inequalities require no limiting
argument.
\end{proof}

\subsection{Three balls and the equality classification}

A normed space has the \emph{three-ball intersection property}, also
called the \emph{3.2 intersection property}, if any three closed balls
that intersect pairwise have a common point. Their radii may differ. The equivalence between this intersection
property and the metric-median condition is Lima's classical
norm-additive decomposition criterion~\cite[Theorem~1.2]{Lima1978};
see also Hansen--Lima~\cite[Theorem~3.6(4)]{HansenLima1981}. We include the
short direction needed here, with arbitrary radii.

\begin{lemma}[From medians to three balls]\label{lem:three-balls}
Suppose every triple of points in a real normed space has a metric
median. Then the space has the three-ball intersection property,
including balls of radius zero.
\end{lemma}

\begin{proof}
Consider pairwise-intersecting balls with centers $x_1,x_2,x_3$ and
radii $r_1,r_2,r_3\ge0$. Let $m$ be a metric median of their centers,
and set $d_i=\|m-x_i\|$. Pairwise intersection and the median property
imply
\[
 d_i+d_j=\|x_i-x_j\|\le r_i+r_j
 \qquad(i\ne j).
\]
If every $d_i\le r_i$, the median is a common point. Otherwise,
after relabeling, $t=d_1-r_1>0$. For $j=2,3$ the displayed inequalities
give $d_j+t\le r_j$, so no other ball excludes $m$. Since
$0<t\le d_1$, the point
\[
 z=m+\frac{t}{d_1}(x_1-m)
\]
is well defined and lies on the segment from $m$ to $x_1$. It satisfies
\[
 \begin{aligned}
 \|z-x_1\|&=d_1-t=r_1,\\
 \|z-x_j\|&\le\|z-m\|+d_j=t+d_j\le r_j\quad(j=2,3).
 \end{aligned}
\]
Thus $z$ lies in all three balls. This includes $r_1=0$, when $z=x_1$.
\end{proof}

The required structure theorem is due to Hansen and Lima. In their
terminology the balls are closed and have arbitrary radii
\cite[Section~1]{HansenLima1981}.

\begin{theorem}[Hansen--Lima]\label{thm:hansen-lima}
A nonzero finite-dimensional real Banach space with the 3.2 intersection
property is linearly isometric to a space obtained from the real line
by finitely many $\ell_1$ and $\ell_\infty$ direct sums.
\end{theorem}

This is \cite[Corollary~7.4]{HansenLima1981}. The two direct-sum norms
are exactly those in \eqref{eq:sum-gauges}, so their unit balls are
Hanner polytopes.

\begin{proposition}[Classification of equality]\label{prop:equality-classification}
If an origin-symmetric convex body $K\subset\R^n$, $n\ge1$, satisfies
$P(K)=4^n/n!$, then it is a linear Hanner body.
\end{proposition}

\begin{proof}
Put $Q=K^\circ$. Proposition~\ref{prop:metric-median} and
Lemma~\ref{lem:three-balls} give the 3.2 intersection property for
$(\R^n,p_Q)$. This is a real Banach space, since every finite-dimensional
normed space is complete. Theorem~\ref{thm:hansen-lima} gives an
invertible linear map $T$ and a Hanner polytope $H$ with $Q=TH$.
By \eqref{eq:linear-polarity} and Lemma~\ref{lem:hanner},
\[
 K=Q^\circ=T^{-\mathsf T}H^\circ
\]
is a linear Hanner body. The converse was proved in
Lemma~\ref{lem:hanner}; together with
Corollary~\ref{cor:symmetric-inequality}, this completes
Theorem~\ref{thm:main}.
\end{proof}
